% !TeX root = ../../Book3.tex
\OptionalSection{二次型的 \texorpdfstring{\(\theta\)}{\AnalysisPdfTheta} 级数}{b3:sec:2405}

把平方和换成一般正定二次型，收敛和生成函数意义仍保留，但参数反演会涉及对偶格，模群也通常缩小。

\subsection{正定二次型}
\label{b3:subsec:240501}

设 \(A\) 为实对称正定矩阵，\(Q(x)=\tfrac12x^TAx\)。正定性给 \(Q(x)\ge c|x|^2\)，因而 \(\sum_{m\in\mathbb Z^d}e^{2\pi\mathrm{i}\tau Q(m)}\) 在上半平面正规收敛。若 \(A\) 为整矩阵且对角线为偶数，则 \(Q(m)\) 为整数，级数对 \(\tau\mapsto\tau+1\) 不变。

\subsection{格 \texorpdfstring{\(\theta\)}{\AnalysisPdfTheta} 级数}
\label{b3:subsec:240502}

对欧氏格 \(L\)，记 \(L^*=\{y:\langle x,y\rangle\in\mathbb Z\ (x\in L)\}\)，基本区域体积为 \(v(L)\)。格 theta 级数为
\(\Theta_L(\tau)=\sum_{x\in L}e^{\pi\mathrm{i}\tau|x|^2}\)。把实高斯的 Fourier 变换按基矩阵换元，Poisson 求和给
\[
\Theta_L(-1/\tau)=\frac{(-\mathrm{i}\tau)^{d/2}}{v(L)}\Theta_{L^*}(\tau).
\]
高维求和公式可先对标准整数格逐个坐标作周期 Fourier 展开，再线性换元；高斯的共同快速衰减保证所有换序。

\subsection{模性}
\label{b3:subsec:240503}

若 \(L=L^*\)，并且每个格向量平方长度为偶整数，则 \(v(L)=1\)，平移与反演都封闭在同一个 theta 级数上。若另有 \(d\) 是八的倍数，则 \((-\mathrm{i})^{d/2}=1\)，两个生成元给出权 \(d/2\) 的全模群变换律；尖点常数项一及指数尾部给尖点全纯性。因此 \(\Theta_L\in M_{d/2}\)。这里把维数条件列为假设，并不在本节证明偶自对偶格的维数必为八的倍数。

\subsection{表示数}
\label{b3:subsec:240504}

若 \(Q\) 在整格上取非负整数值，则 Fourier 系数等于表示数 \(r_Q(n)=\#\{m:Q(m)=n\}\)。正定性保证每层有限。对非自对偶格，反演把它送到对偶格或若干陪集级数的组合，通常得到子群上的模形式或向量值变换律；不能从一个高斯公式直接宣称所有整二次型都给全模群标量模形式。\appref{b3:app:E}将说明有限层级的相应语言。

\ProseParagraphBreak

高维 Gaussian 公式中的体积因子可在矩阵形式下直接看到。若 \(A\) 正定，选正交坐标将它对角化，再逐坐标使用一维 Gaussian 变换，得到
\[
 \widehat{e^{-\pi t x^TAx}}(\xi)
 =t^{-d/2}(\det A)^{-1/2}e^{-\pi\xi^TA^{-1}\xi/t}.
\]
正定性确保特征值全正、Gaussian 可积；逆矩阵出现在频率的二次型中，行列式来自实变量换元。因此反演自然交换原格与对偶格，而不会对任意正定整矩阵保持同一个级数。若二次型只是半正定，例如 \(Q(x,y)=x^2\)，每个固定表示层中 \(y\) 可任意变化，系数已经无穷；若不定，例如 \(x^2-y^2\)，零层也无穷。正定性不仅用于估计，还保证生成函数本身的计数意义。
\ProseParagraphBreak
偶自对偶条件有一个具体的八维实例。令
\[
 D_8=\{x\in\mathbb Z^8:\textstyle\sum_jx_j\text{ 为偶数}\},\qquad
 L=D_8\cup\left(D_8+\left(\tfrac12,\ldots,\tfrac12\right)\right).
\]
两次半整数平移回到 \(D_8\)，故并集是加法群。整数部分的平方长度模二等于坐标和，因而偶；半整数部分 \(x+v\) 的平方长度为 \(|x|^2+\sum x_j+2\)，也为偶整数。由极化恒等式，任意两格向量的内积为整数，所以 \(L\subset L^*\)。\(D_8\) 在整数格中指标为二，体积二；\(L\) 又含它为指标二子格，故体积一。对偶格体积为倒数一，包含关系与相同体积迫使 \(L=L^*\)。这给出本节模性结论所需的全部假设。
\ProseParagraphBreak
还可独立数出其最短非零向量。整数部分长度平方为二的向量恰有两个坐标为 \(\pm1\)，共 \(\binom82\cdot4=112\) 个。半整数部分若有坐标绝对值至少 \(3/2\)，平方长度已至少四；所以长度平方二时八个坐标全为 \(\pm1/2\)。属于指定陪集要求负号数为偶数，恰有 \(2^7=128\) 个。总数为240，与权四 Eisenstein 级数的一次系数相符。这个几何核算检验了长度平方 \(2n\) 与模变量 \(e^{2\pi\mathrm{i}\tau}\) 中 \(q^n\) 的对应，也给抽象的偶自对偶条件一个可以逐坐标检查的模型。
